
\documentclass[12pt,a4paper]{article}

\usepackage[utf8]{inputenc}
\usepackage[T1]{fontenc}
\usepackage[brazil]{babel}
\usepackage{amsmath}
\usepackage{siunitx}
\usepackage{booktabs}
\usepackage{graphicx}
\usepackage[margin=2.5cm]{geometry}

\graphicspath{{figs/}}

\begin{document}

\subsection*{Caso 3S}
\label{sec:res-caso3s}

Removida toda a compressão axial, a única carga mecânica adicionada é a pressão focal da artéria oftálmica como carga lateral prescrita ($p_c = \SI{9034}{\pascal}$, sentido $+x$, em $z \approx \SI{22.5}{\milli\meter}$); permanece apenas a PIC basal de \SI{1333}{\pascal} aplicada na face interna da dura (\texttt{dura\_inner}, Tab.~\ref{tab:casos-setup}). O solver de Riks satura monotonicamente em $\lambda = 1{,}0$ com span axial residual de apenas $\approx \SI{36}{\micro\meter}$, confirmando que a deformação provém apenas da carga lateral, sem componente de flambagem axial global. Sob carga plena a artéria produz uma indentação lateral localizada de sentido único, concentrada na parede da bainha sob o contato (excursão da face externa da dura $\approx \SI{606}{\micro\meter}$) e decaindo em direção ao eixo do nervo ($|U_x|_{\mathrm{eixo}} \approx \SI{333}{\micro\meter}$, resultante lateral $\approx \SI{71.7}{\milli\newton}$), e não um arco global nem a ondulação antissimétrica em ``S''.

\end{document}
